\color{unidarkgreen}{\subsection{Токовая защита обратной последовательности}}
\color{black}

\begin{enumerate}

\item
Токовая защита обратной последовательности (ТЗОП) предназначена для резервирования действия основных быстродействующих защит при внутренних несимметричных коротких замыканиях в БСК. Защита выполнена одноступенчатой на базе модуля \textbf{«ТЗОП»}. Выводится защита накладкой \textbf{«ТЗОП ВВ»}. Измерительные цепи защиты включены на токи фаз А, В, С со стороны высоковольтных вводов БСК. Защита реагирует на повышение тока обратной последовательности, определяемого по формуле:

\begin{equation}
\label{eq:I2}
\dot{I}_{2} = \frac{\dot{I}_{\textsl{A}} + a^2 \cdot \dot{I}_{\textsl{B}} + a \cdot \dot{I}_{\textsl{C}}}{3}
\end{equation}

\item
При необходимости защита может выполнена с блокировкой по второй гармонике (отношению уровня второй гармоники к первой). Данная блокировка реализована на базе модуля \textbf{«ОВ БНТ ВВ»}. Выводится данная блокировка накладкой \textbf{«Режим БНТ ТЗНП-1 ВВ»}.

\item
Защита имеет регулируемые значения по току и времени срабатывания, задаваемые параметрами \textbf{«Iср ТЗОП ВВ»} и \textbf{«Тср ТЗОП ВВ»} соответственно. 

\item 
Уставки ТЗОП приведены в таблице \ref{app_set:tzop}. Алгоритм работы ТЗОП приведен на рисунке \ref{pict:TZOP}.

\medskip
\begin{figure}[H]
\centering
\includegraphics[width=1\textwidth,height=1\textheight,keepaspectratio]{\fbpath/09.01 БСК/271.0291 ТЗОП/img/TZOP.pdf}
\captionof{figure}{Функциональная схема логики ТЗОП}\label{pict:TZOP}
\end{figure}


\end{enumerate} 